%% ============================================================
%% Chapter 8: TARTAN: Recursive Multiscale Tiling
%% ============================================================

\chapter{TARTAN: Recursive Multiscale Tiling}
\label{ch:ch08-tartan-tiling}
\section{Introduction}

Neither Chapter~\ref{ch:ch06-distinction-holonomy}'s continuous fibers nor Chapter~\ref{ch:ch07-spherepop}'s discrete events, by themselves, supply the geometry connecting them: a principled account of which regions of a continuous field admit one coherent local continuation law, and how that account changes under refinement. TARTAN, Trajectory-Aware Recursive Tiling with Annotated Noise, provides this multiscale structure, recursively partitioning the history space into coherence tiles, each carrying a local field state, entropy, noise annotation, scale, and trajectory estimate.

The methodological precedent is Turing's, in two distinct respects: the theory of computation shows how complex behavior can be generated by composition of elementary operations, and the reaction-diffusion account of morphogenesis shows how local interactions can generate spatial pattern without a central pattern-maker. TARTAN combines these only at the level of formal organization, complex global structure analyzed through local rules, spatial coupling, and the repeated differentiation of initially less differentiated regions, and neither claims equivalence to reaction-diffusion dynamics nor that every tiled system reproduces biological morphogenesis. Spherepop then governs changes to the tiling:
\[
\Pop \leftrightarrow \text{selection and refinement}, \quad \Refuse \leftrightarrow \text{exclusion of an inadmissible adjacency},
\]
\[
\Bind \leftrightarrow \text{coupling between tiles}, \quad \Collapse \leftrightarrow \text{projection or coarse-graining}.
\]
Distinction holonomy measures what survives transport through the resulting recursively changing tile system.

\section{Tiles and Coherence}

Let the RSVP field state be $\Psi = (\Phi, \mathbf{v}, S)$ over a domain $M$. At scale $k$, a tiling $\mathcal{T}_k = \{T_i^{(k)}\}_{i \in I_k}$ covers $M$, with each tile
\[
T_i^{(k)} = \bigl(U_i^{(k)}, \Psi_i^{(k)}, A_i^{(k)}, \eta_i^{(k)}, \tau_i^{(k)}, H_i^{(k)}\bigr),
\]
recording support, local RSVP state, local continuation connection, annotated noise, predicted trajectory, and local append-only history. Recursive tiling allows a tile to be replaced by children, $T_i^{(k)} \rightsquigarrow \{T_{i1}^{(k+1)}, \ldots, T_{im}^{(k+1)}\}$ with $U_i^{(k)} = \bigcup_a U_{ia}^{(k+1)}$, giving a refinement structure $\mathcal{T}_0 \preceq \mathcal{T}_1 \preceq \mathcal{T}_2 \preceq \cdots$: not one base geometry, but a recursively changing family of bases.

A tile should be formed where field dynamics admit a sufficiently coherent local continuation model. With a coherence functional
\[
\kappa = \alpha_\Phi \|\nabla \Phi\|^2 + \alpha_v \|\nabla \mathbf{v}\|^2 + \alpha_S \|\nabla S\|^2 + \alpha_F \|F\|^2 + \alpha_\eta \operatorname{Var}(\eta),
\]
a tile is a connected region $\{x \in M : \kappa(x,t) \leq \varepsilon_i\}$: an improvement over defining coherence from entropy gradient alone, since a region can have small $|\nabla S|$ either because it is genuinely coherent or because it has been destructively homogenized, and the curvature and noise terms help distinguish these. A coherence tile is a region over which transport is approximated by one local connection $\nabla_i = d + A_i$ with curvature $F_i = dA_i + A_i \wedge A_i$; low curvature means nearby alternative paths within the tile transport distinctions approximately equivalently, not that the tile is globally memoryless.

\section{The Four Primitives as Tiling Operations}

$\Pop$ becomes trajectory-aware refinement: if a tile $T$ contains several distinguishable continuation regimes, $\Pop_\omega(T) = (T_\omega, T_{\neg\omega})$ with $T_\omega = \{x \in T : \omega \text{ selected}\}$, and the selection is committed to history. Pop therefore operates at two levels at once, choosing one continuation, and introducing a new tile boundary that is itself a realized distinction: commitment that becomes geometry. Trajectory awareness means refinement can depend on predicted divergence, $d(\tau_\omega(t+\Delta t), \tau_{\neg\omega}(t+\Delta t)) > \delta$, so states locally similar now may be separated because their admissible futures diverge.

$\Refuse$ becomes prohibited gluing: for adjacent tiles $T_i, T_j$ with a proposed gluing map $g_{ij} : E_i|_{U_i \cap U_j} \to E_j|_{U_i \cap U_j}$, $\Refuse(g_{ij})$ removes $g_{ij}$ from the admissible edges of the tile-adjacency graph $\mathcal{G}$ while the refusal itself remains in history. Refuse denies a proposed overlap identification, records that a proposed output fails an interface contract, or removes a candidate transport from the admissible connection, three equivalent readings in sheaf, type-theoretic, and gauge language respectively.

$\Bind$ becomes intertile coupling, $\Bind(T_i,T_j) : T_i \to T_j$, represented by an operator $B_{ij} : E_i \to E_j$ or a coupled-field action term $\sum_{i,j} \int_{U_i \cap U_j} \|s_j - B_{ij} s_i\|^2\, dV$; the local dynamics becomes $\partial_t \Psi_j = \mathcal{F}_j[\Psi_j] + \lambda_{ij} \mathcal{B}_{ij}[\Psi_i,\Psi_j]$. Because binding changes both the tile graph and the local dynamics, it is again the natural source of curvature: $\Pop_i \circ \Bind_{ij} \neq \Bind_{ij} \circ \Pop_i$ in general, since the two orders generate different descendant couplings.

$\Collapse$ becomes recursive coarse-graining: given child tiles $\{T_{i1}^{(k+1)}, \ldots, T_{im}^{(k+1)}\}$, a certificate $c$ defines $\pi_c : \coprod_a T_{ia}^{(k+1)} \twoheadrightarrow T_i^{(k)}$, admissible only when constitutive structure descends through the quotient, a coarse binding $\overline{B}$ with $\pi_c \circ B = \overline{B} \circ \pi_c$, and an effective coarse connection $A_i^{(k)} = \mathcal{R}_c(A_{i1}^{(k+1)}, \ldots, A_{im}^{(k+1)})$ for a renormalization operator $\mathcal{R}_c$. Collapse is destructive exactly when $\pi_c(x) = \pi_c(y)$ but $\Hol(x) \not\sim \Hol(y)$ for distinct child histories: coarse-graining certified by continuation equivalence.

\section{Annotated Noise and Induced Curvature}

Each tile carries a noise annotation $\eta_i = (\mu_i, \Sigma_i, \ell_i, \chi_i)$, local bias, covariance, correlation scale, provenance, rather than one undifferentiated random term. The local connection becomes stochastic, $A_i = \overline{A}_i + \Xi_i$ with $\mathbb{E}[\Xi_i] = \mu_i$, and the curvature is correspondingly random, with expectation
\[
\mathbb{E}[F_i] = d\overline{A}_i + \overline{A}_i \wedge \overline{A}_i + \mathbb{E}[\Xi_i \wedge \Xi_i] + \text{cross terms}.
\]
Consequently $\mathbb{E}[\Xi_i] = 0$ does not imply $\mathbb{E}[F_i] = F(\overline{A}_i)$: zero-mean connection noise can generate nonzero effective curvature. Correlated noise need not merely obscure an underlying continuation geometry; it can create systematic holonomy.

\section{Three Scales of Curvature}

Curvature separates into three levels. Within a tile, differential curvature is $F_i = dA_i + A_i \wedge A_i$ as usual. Across a boundary, transition curvature measures incompatibility between local connections: compatible transport requires $A_j = g_{ij}^{-1} A_i g_{ij} + g_{ij}^{-1} dg_{ij}$, and the boundary defect
\[
K_{ij} = A_j - \bigl( g_{ij}^{-1} A_i g_{ij} + g_{ij}^{-1} dg_{ij} \bigr)
\]
signals that the two tiles cannot be treated as local descriptions of one smoothly transported structure when $K_{ij} \neq 0$. Across scales, recursive curvature measures the failure of refinement and transport to commute: for a coarse-graining map $r_{k+1,k}$, scale compatibility requires $r_{k+1,k} \circ \mathsf{T}^{(k+1)}_\gamma = \mathsf{T}^{(k)}_\gamma \circ r_{k+1,k}$, and the defect
\[
K_\gamma^{(k+1,k)} = r_{k+1,k} \circ \mathsf{T}^{(k+1)}_\gamma - \mathsf{T}^{(k)}_\gamma \circ r_{k+1,k}
\]
being nonzero means a history visible at the fine scale is misrepresented after coarse-graining: the system can appear flat macroscopically while retaining hidden microscopic holonomy.

\section{Recursive Holonomy and Multiscale Persistence}

A trajectory through TARTAN moves among tiles and scales, $\gamma : T_{i_0}^{(k_0)} \to T_{i_1}^{(k_1)} \to \cdots \to T_{i_n}^{(k_n)}$, each step an operator $U_m : E_{i_{m-1}}^{(k_{m-1})} \to E_{i_m}^{(k_m)}$ that may be a continuous transport, a Spherepop event, an intertile binding, a refinement, or a quotient projection. The TARTAN holonomy is the path-ordered composite $\Hol_\gamma^{\mathrm{TARTAN}} = U_n \circ \cdots \circ U_1$, generally an endomorphism rather than a group element. Even for a loop returning to the same coarse tile, $T_{i_n}^{(k_n)} = T_{i_0}^{(k_0)}$, one may have $\Hol_\gamma^{\mathrm{TARTAN}} \neq I$: the residual transformation records not only where the trajectory traveled but which distinctions were selected, excluded, coupled, or collapsed, and at which scales. TARTAN memory is holonomy through space, history, and scale.

At each scale $k$, let $\mathfrak{I}_k$ denote the continuation identity represented by $\mathcal{T}_k$, with coarse-graining maps $r_{k+1,k} : \mathfrak{I}_{k+1} \to \mathfrak{I}_k$. A scale-coherent persistent entity is a compatible family $(\mathfrak{I}_0, \mathfrak{I}_1, \ldots)$ with $r_{k+1,k}(\mathfrak{I}_{k+1}) = \mathfrak{I}_k$, and the complete entity is the inverse limit
\[
\mathfrak{I} = \varprojlim_k \mathfrak{I}_k.
\]
An object need not exist at one privileged resolution; it exists as compatibility among its representations across resolutions. A failure confined to one scale need not destroy the entity; ontological failure in the TARTAN sense occurs only when no compatible family remains, $\varprojlim_k \mathfrak{I}_k = \varnothing$, stronger than disappearance from any single tiling.

\section{TARTAN Repair}

Damage may occur within a tile, at a boundary, or between scales, so repair has three components, $R = (R_{\mathrm{local}}, R_{\mathrm{glue}}, R_{\mathrm{scale}})$, restoring intratile transport, gluing compatibility, and cross-scale agreement respectively, with combined functional
\[
\mathcal{L}_{\mathrm{TARTAN}}[R] = \sum_i \int_{U_i} \|F_i^R - F_i^{\mathrm{ref}}\|^2\, dV \;+\; \lambda_b \sum_{i,j} \|K_{ij}^R\|^2 \;+\; \lambda_s \sum_k \|K^{(k+1,k),R}\|^2 \;+\; \lambda_R \|R\|^2,
\]
minimized subject to preserving the constitutive holonomy observables. The distinctive TARTAN possibility is retiling: if a damaged tile can no longer sustain one coherent continuation law, repair need not restore its original interior. Spherepop can $\Pop$ it into several tiles, $T_i \to \{T_{i1}, T_{i2}\}$, $\Bind$ their viable dependencies, $\Refuse$ destructive crossings, and $\Collapse$ distinctions irrelevant at the repaired scale: repair as adaptive repartitioning of the future, and the operation Chapter~\ref{ch:ch07-spherepop}'s double-bind discussion invokes by name as the only route out of a trapped level.

\section{The Combined Architecture}

The complete architecture reads
\[
\text{RSVP} \to \text{continuous field dynamics},
\]
\[
\text{TARTAN} \to \text{trajectory-aware multiscale decomposition},
\]
\[
\text{Spherepop} \to \text{event grammar transforming that decomposition},
\]
\[
\text{DHT} \to \text{transport, curvature, holonomy, persistence}.
\]
RSVP supplies the evolving substrate $\Psi = (\Phi,\mathbf{v},S)$; TARTAN identifies the regions and scales over which local continuation laws are coherent; Spherepop changes those regions through Selection, Exclusion, Coupling, and Projection; and Chapter~\ref{ch:ch06-distinction-holonomy}'s apparatus measures the historical effect of transporting constitutive distinctions through the resulting structure. Their noncommutativity generates curvature; their path-ordered composition generates holonomy; TARTAN makes both recursive and scale-dependent. The resulting picture departs from a single-tile idealization in one significant respect: an entity is not represented by one distinction fiber over one base manifold. It is represented by the compatibility of its local continuation structures across a recursively changing tiling: a persistent entity is a multiscale descent object whose constitutive holonomy survives admissible retiling.

\section{An Independent Instance: Contextuality Without a Global Section}

TARTAN's boundary defect $K_{ij}$ and its inverse-limit criterion for persistence, $\mathfrak{I} = \varprojlim_k \mathfrak{I}_k \neq \varnothing$, are both instances of a much older mathematical question: when can locally consistent data be glued into one globally consistent object? A recent paper by Buchanan, developing a general categorical framework for context-indexed observation, arrives at a version of this question from an entirely different direction, quantum foundations, and is worth reading here as an independent case study, not because it was written with TARTAN in mind, but because the convergence itself is informative about how general the underlying obstruction pattern is.

Buchanan's framework organizes contexts into a category $\mathbf{Ctx}$, with observation modeled as a bifunctor $\Phi : \mathbf{Agent} \times \mathbf{Ctx}^{\mathrm{op}} \to \mathbf{Cat}$ sending each agent-context pair to its category of evaluations, and a Grothendieck construction assembling these into a single fibered category $\int \Phi \to \mathbf{Agent} \times \mathbf{Ctx}$. Two evaluations by different agents under different contexts are \emph{synonymous} when they agree despite the difference in context; the paper shows, under an explicit stated hypothesis about the existence of a common refinement, that synonymy of two localized evaluations does not by itself entail that the restriction map producing them is an isomorphism, only that it identifies those two particular elements. This is exactly TARTAN's caution around $K_{ij}$: two adjacent tiles agreeing on their overlap does not certify that the gluing map $g_{ij}$ is globally coherent, only that this particular comparison happened to succeed.

The paper's established core result, developed via the sheaf-theoretic account of Kochen--Specker contextuality, represents non-contextuality as the existence of a global section, a single evaluation consistent with every local agent-context restriction simultaneously, and represents the established quantum-mechanical failure of non-contextuality as the non-existence of such a global section for suitable measurement scenarios: local evaluations that are individually consistent on every overlap can nonetheless admit no globally coherent assignment. This is, formally, the Kochen--Specker analogue of TARTAN's own persistence criterion stated in the negative: $\varprojlim_k \mathfrak{I}_k = \varnothing$ despite every individual $\mathfrak{I}_k$ being well-defined and every finite overlap being locally consistent. Where TARTAN attributes such a failure to accumulated curvature, boundary defects $K_{ij}$ and cross-scale defects $K_\gamma^{(k+1,k)}$ that do not cancel, the Kochen--Specker result exhibits a concrete physical system for which no choice of local data, however carefully arranged, can make those defects vanish everywhere at once: the obstruction is not a failure of engineering but a structural feature of the underlying measurement scenario.

This reading is offered with the same care the reproducibility discussion of Chapter~\ref{ch:ch04-paracosm-trap} insisted on: it is a structural analogy recognized after the fact, not a claim that Buchanan's paper was doing TARTAN theory under another name, and not a claim that the vocabularies straightforwardly translate. Buchanan's own paper is explicit that several of its further proposed connections, to homological extension classes, to Kan-extension reconstruction, and to deformation-quantization cocycles, are open comparison problems rather than proved results, a discipline about the boundary between established and speculative content worth naming here too. What the established core of the two frameworks do genuinely share, independently arrived at, is the same basic shape: local coherence is cheap, global coherence is not automatic, and the gap between them is where a system's persistence, or a measurement scenario's classicality, actually lives or fails to.
